\documentclass[11pt,a4paper]{article}
\usepackage[margin=25mm]{geometry}
\usepackage{amsmath,amssymb,amsthm,hyperref}
\setlength{\emergencystretch}{2em}
\newtheorem{lemma}{Lemma}
\newtheorem{corollary}[lemma]{Corollary}
\newcommand{\gen}[1]{\langle#1\rangle}
\title{A fixed-support exit from an internal pair equality}
\author{Human Compression project: L044 continuation}
\date{}
\begin{document}\maketitle
Retain the complete hypotheses of Theorem 4 of
\emph{Ambient transport and internal pair deficits}, including
$k\notin F$, intransitivity of the full local $K=G_{>k}$ on the
actual $H=G_{\ge k}$-orbit $O$, and equality $|F|=q-1$ with
$|O|=2q$, $q\ge3$. Thus $K|_O$ is the diagonal $S_q$ on two
natural rows, and every tail mark outside $F\cup\{k\}$ fixes $O$.
Let $a=\min F$ and $x=\max F=a+q-2$; here $a$ denotes a label,
not the crossing vertex used elsewhere. Put $P=G_{<k}$.
Suppose also that a specified set $C\subseteq\operatorname{Fix}(K)$
has size $h\ge1$. The ambient group $G$ is transitive.

\begin{lemma}[A third fixed row]\label{rows}
Every point of the union $O^P=\bigcup_{p\in P}O^p$ is fixed by
$G_{>x}$. This union is a disjoint union of $K$-orbits of size $q$,
is disjoint from $C$, and has at least $3q$ points. Consequently
\[
                |\operatorname{Fix}(G_{>x})|\ge 3q+h.    \tag{1}
\]
\end{lemma}
\begin{proof}
The two rows of $O$ are actual $K$-orbits of size $q$.
The prefix $P$ centralizes $K$, so each of their translates is
again a $K$-orbit of size $q$. Distinct orbits are disjoint and
none contains a point of $C$, all of whose points are fixed by $K$.
If $O^P=O$, then $P$ preserves $O$; so does $H$, by its definition.
Since $G=\gen{P,H}$, transitivity would give $O=\Omega$.
This is impossible because $C$ is nonempty and disjoint from $O$.
Thus at least one further $q$-point $K$-orbit occurs in $O^P$.

Every mark after $x$ lies outside $F\cup\{k\}$ and hence fixes
$O$ pointwise. Also $G_{>x}\le K$, so it commutes with $P$.
For $v\in O$, $p\in P$ and $g\in G_{>x}$, one has
$(v^p)^g=(v^g)^p=v^p$. Thus $G_{>x}$ fixes $O^P$ pointwise.
It fixes $C$ as a subgroup of $K$. The disjointness already proved
gives (1). This argument does not require $K$ to be transitive
on any larger set or identify two distinct row orbits.
\end{proof}

\begin{corollary}[The exact source exit]\label{exit}
Put $X=\Omega\setminus\operatorname{Fix}(G_{>x})$ and
$n=|\Omega|$. If
\[
                         2a\le q+h+3,                  \tag{2}
\]
then
\[
                         x\le\frac{n-|X|-1}{2}.        \tag{3}
\]
In particular this holds when the internal path starts at $k+1$
in either source head configuration $(k,h)=(2,2)$ or $(5,7)$.
This is alternative (a) of source Proposition 5.31.
\end{corollary}
\begin{proof}
Since $x=a+q-2$, condition (2) says
$2x+1=2a+2q-3\le3q+h$. Lemma~\ref{rows} gives
$2x+1\le n-|X|$, which is (3). If $a=k+1$, condition (2)
reduces to $q\ge1$ for $(k,h)=(2,2)$ and to $q\ge2$ for
$(k,h)=(5,7)$; both follow from $q\ge3$.
There is no assumption $x<r-1$: if $x=r-1$, the fixed set in
(1) is all of $\Omega$ and the same inequality remains valid.
\end{proof}

\paragraph{Source and scope.}
The target is Cameron--Fernandes--Leemans--Mixer,
\emph{Highest rank of a polytope for $A_n$}, arXiv:1605.09173v1,
Proposition 5.31(a) and its internal-orbit cases, physical
pages 38 and 42--43. Its two head configurations supply two or
seven specified $K$-fixed points; an application must identify
these actual points. The fixed-support target is borrowed; the
third-row proof avoids discarding the even flip pattern. The full
internal-pair proof and its ambient-eligibility hypotheses remain
paid. A later path is covered only under (2). Other deficits,
tables and full L044 remain pending; no computation, formal theorem
or question credit is added.
\end{document}
